\documentclass[11pt]{article}
\usepackage[a4paper,margin=27mm]{geometry}
\usepackage{amsmath,amssymb,amsthm,xcolor,booktabs,array}
\newcommand{\PROVED}{\textcolor{green!40!black}{\textbf{PROVED}}}
\newcommand{\OPEN}{\textcolor{red!70!black}{\textbf{OPEN}}}
\newcommand{\AUDIT}{\textcolor{orange!70!black}{\textbf{AUDIT}}}
\newtheorem{theorem}{Theorem}
\newtheorem{corollary}{Corollary}
\title{Exact newborn arithmetic cancellation and Gamma coercivity (v118)}
\date{13 August 2026}

\begin{document}
\maketitle

\section{A support-only lower bound for the Gamma energy}

Use the Fourier convention
\[
 \widehat f(\tau)=\int_{\mathbb R}f(t)e^{i\tau t}\,dt,
 \qquad
 \|f\|_2^2=\int_{\mathbb R}|\widehat f(\tau)|^2\frac{d\tau}{2\pi}.
\]
At the critical normalization the positive Gamma multiplier is
\begin{equation}
 g_\Gamma(\tau)=\sum_{j\ge0}
 \frac{2\tau^2}{a_j(a_j^2+\tau^2)},
 \qquad a_j=2j+\frac12.                                     \tag{1}
\end{equation}

\begin{theorem}[Logarithmic uncertainty bound --- \PROVED]
If $f\in L^2(\mathbb R)$ is supported on a measurable set of measure
$m\le\pi/2$, then
\begin{equation}
 \boxed{
 \int_{\mathbb R}g_\Gamma(\tau)|\widehat f(\tau)|^2
                 \frac{d\tau}{2\pi}
 \ge \frac38\log\!\left(\frac{\pi}{2m}\right)\|f\|_2^2.} \tag{2}
\end{equation}
\end{theorem}

\begin{proof}
For $\tau\ge\tfrac12$, retain in (1) all $j$ such that $a_j\le\tau$.
Each retained summand is at least $1/a_j$.  If
$J=\lfloor(\tau-\frac12)/2\rfloor$, then
\begin{align*}
 g_\Gamma(\tau)
 &\ge\sum_{j=0}^J\frac1{2j+1/2}
 \ge\int_0^{J+1}\frac{dx}{2x+1/2}\\
 &=\frac12\log(4J+5)
 \ge\frac12\log(2\tau).                                   \tag{3}
\end{align*}
The same bound holds with $|\tau|$ by evenness.

Cauchy--Schwarz on the support gives
$|\widehat f(\tau)|^2\le m\|f\|_2^2$.  Put
$R=\pi/(4m)\ge1/2$.  Then
\begin{equation}
 \int_{-R}^R|\widehat f(\tau)|^2\frac{d\tau}{2\pi}
 \le\frac{Rm}{\pi}\|f\|_2^2
 =\frac14\|f\|_2^2.                                       \tag{4}
\end{equation}
Thus at least three quarters of the Fourier mass lies in $|\tau|\ge R$.
Equations (3)--(4), positivity of $g_\Gamma$, and monotonicity of every
summand of (1) in $|\tau|$ give
\[
 \langle G_\Gamma f,f\rangle
 \ge\frac34\,g_\Gamma(R)\|f\|_2^2
 \ge\frac38\log(2R)\|f\|_2^2,
\]
and $2R=\pi/(2m)$.  This proves (2).
\end{proof}

\section{The newborn arithmetic block}

Put
\begin{equation}
 T_N=\frac12\log N,
 \qquad \delta_N=T_{N+1}-T_N,
 \qquad
 E_N=(-T_{N+1},-T_N)\cup(T_N,T_{N+1}).                       \tag{7}
\end{equation}
Let $S_a$ denote translation after zero extension.

\begin{theorem}[Exact vanishing of every newborn self-contact ---
\PROVED]
For every integer $n$ with $2\le n\le N$ and every
$f\in L^2(E_N)$,
\begin{equation}
                    \langle(S_{\log n}+S_{-\log n})f,f\rangle=0.
                                                                    \tag{8}
\end{equation}
Consequently the complete arithmetic part of the localized row-(d)
operator has zero newborn/newborn block.
\end{theorem}

\begin{proof}
Each component of $E_N$ has length $\delta_N$.  Since
$\log n\ge\log2>\delta_N$, a translated component cannot overlap itself.
An overlap between opposite components would require simultaneously
\[
 2T_N<\log n<2T_{N+1},
\]
which is exactly $N<n<N+1$.  No integer satisfies this.  Endpoint contacts
have measure zero.  Hence every correlation in (8) vanishes.
\end{proof}

Let
\begin{equation}
 A_{T_{N+1}}=G_{\Gamma,T_{N+1}}-m_0I-
 \sum_{2\le n<N+1}\frac{\Lambda(n)}{\sqrt n}
        (S_{\log n}+S_{-\log n}).                            \tag{9}
\end{equation}

\begin{corollary}[Exact raw newborn diagonal --- \PROVED]
The $E_N$ diagonal block of (9) is
\begin{equation}
                         A_{EE}=G_{\Gamma,EE}-m_0I.           \tag{10}
\end{equation}
Since $|E_N|=2\delta_N$, (2) gives
\begin{equation}
 A_{EE}\ge
 \left[
  \frac38\log\!\left(\frac{\pi}{4\delta_N}\right)
  -m_0\right]I.                                             \tag{11}
\end{equation}
In particular the raw newborn diagonal is strictly positive for all
sufficiently large $N$ and grows at least as
$\frac38\log N+O(1)$.
\end{corollary}

\begin{proof}
Equation (10) follows from Theorem 2 term by term.  Apply (2) with
$m=2\delta_N$.  Finally
$\delta_N=\frac1{2N}+O(N^{-2})$.
\end{proof}

\section{Exact norm of the arithmetic old/new cross}

Let $P_{0,N}$ be restriction to the old interval $I_N=(-T_N,T_N)$ and
define
\begin{equation}
 B_N^{\rm ar}:=-P_{0,N}\sum_{2\le n\le N}
       \frac{\Lambda(n)}{\sqrt n}(S_{\log n}+S_{-\log n})
       \big|_{L^2(E_N)}.                                    \tag{12}
\end{equation}

\begin{theorem}[Orthogonal arithmetic cross synthesis --- \PROVED]
One has the exact norm identity
\begin{equation}
 \boxed{
 \|B_N^{\rm ar}f\|_2^2
 =\left(\sum_{2\le n\le N}\frac{\Lambda(n)^2}{n}\right)
   \|f\|_2^2,\qquad f\in L^2(E_N).}                         \tag{13}
\end{equation}
In particular
\begin{equation}
 \|B_N^{\rm ar}\|^2
 =\frac12(\log N)^2+o((\log N)^2).                          \tag{14}
\end{equation}
\end{theorem}

\begin{proof}
Split $f=f_-\oplus f_+$ on the left and right components of $E_N$.
After restriction to $I_N$, only $S_{-\log n}f_-$ and
$S_{\log n}f_+$ survive.  Their supports are precisely the intervals
$L_n$ and $R_n$ of v111.  The no-overlap theorem there says that all these
$2N-2$ ranges are pairwise orthogonal.  Every surviving translation is an
isometry, so Pythagoras gives (13).

Only prime powers contribute to the sum.  The primes give, by partial
summation and the prime number theorem,
\[
 \sum_{p\le N}\frac{(\log p)^2}{p}
 =\frac12(\log N)^2+o((\log N)^2).
\]
The powers $p^k$, $k\ge2$, contribute $O(1)$ because
$\sum_p(\log p)^2/[p(p-1)]<\infty$.  This proves (14).
\end{proof}

Explicitly,
\[
 \sum_p\sum_{k\ge2}\frac{(\log p)^2}{p^k}
 =\sum_p\frac{(\log p)^2}{p(p-1)}<\infty.                   \tag{15}
\]

This result identifies the remaining large-cell issue sharply.  It is not
the sign of the raw newborn diagonal.  It is the return through the old
space after the moving-Tate graph of v117:
\begin{equation}
 B_N^*A_{00}^\dagger B_N,
 \qquad
 B_N=P_0(A_{0E}-A_{00}C_N)M_N^{-1/2}.                        \tag{16}
\end{equation}
The prime part of $A_{0E}$ has pairwise orthogonal delay ranges, while the
Gamma part is the Carleman-bounded cross of v112 and
$\|C_N\|=O(N^{-1/2})$ by v117.  Turning these three facts into a sharp
bound strictly below (11) is the physical return estimate; none of
(8)--(11) assumes it.

\begin{center}
\begin{tabular}{>{\raggedright\arraybackslash}p{.55\linewidth}
                >{\centering\arraybackslash}p{.15\linewidth}
                >{\raggedright\arraybackslash}p{.20\linewidth}}
\toprule
Statement & Status & Location \\
\midrule
Support-only Gamma lower bound & \PROVED & (2) \\
All newborn arithmetic self-correlations vanish & \PROVED & (8) \\
Raw newborn block is exactly Gamma minus $m_0$ & \PROVED & (10) \\
Large-$N$ raw newborn positivity & \PROVED & (11) \\
Exact arithmetic old/new cross norm & \PROVED & (13)--(15) \\
Sharp old-space return domination & \OPEN & final display \\
Condition $G$, row (d) & \OPEN & requires (16) in every cell \\
\bottomrule
\end{tabular}
\end{center}

\end{document}
